\documentclass[10pt,a4paper]{amsart}
\usepackage[margin=19mm]{geometry}
\usepackage{amsmath,amssymb,mathtools}
\usepackage[scr=rsfs,cal=euler]{mathalfa}
\usepackage{xfrac}
\usepackage[unicode,hidelinks,bookmarks=false]{hyperref}
\newcommand{\ie}{i.e.\ }
\newcommand{\cf}{cf.\ }
\newcommand{\resp}{respectively\ }
\newcommand{\Subsection}{\subsection}
\providecommand{\nobreakdash}{\nobreak\hbox{-}}
\setlength{\emergencystretch}{2em}
\multlinegap0pt
\usepackage{bm}
\usepackage[normalem]{ulem}
\usepackage{amssymb}
\usepackage{mathtools}
\usepackage{xcolor}
\newtheorem{theorem}{Theorem}
\newcommand{\codenumber}[1]{#1}
\newcommand{\conj}[1]{\overline{#1}}
\DeclareMathOperator{\real}{Re}
\newcommand{\resultnumber}[1]{#1}
\title{Non-uniqueness for the Complex Ginzburg--Landau Equation}
\author{Joel Dahne, Vladimír Šverák}
\date{Source: arXiv:2609.08985v1 (CC BY 4.0)}
\begin{document}
\maketitle

\noindent\emph{Source and licence.} Non-uniqueness for the Complex Ginzburg--Landau Equation. Version arXiv:2609.08985v1 (CC BY 4.0), distributed by arXiv under the Creative Commons Attribution 4.0 International licence, http://creativecommons.org/licenses/by/4.0/, as recorded in the arXiv licence metadata for this exact version and shown on the abstract page https://arxiv.org/abs/2609.08985v1. Full licence notice: \texttt{licenses/arxiv-2609.08985-CC-BY-4.0.txt}.\par\smallskip

\noindent\textit{Editorial scope.} Counted unit: the article's theorem thm:main (source0.tex lines 1484--1503). The article's complete original proof of it is delivered with it (source0.tex lines 1505--1661, one \texttt{proof} environment of 157 lines). No mathematics is written, repaired or re-derived: the statement and the proof are the article's own lines, in the article's own notation, and no retained line is substituted or re-flowed.  Retained uncounted as the source-local context the proof needs: lines 141--146; lines 261--263; lines 366--368; lines 1152--1165; lines 1368--1383; lines 1475--1480; lines 1721--1726; lines 1952--1969.  Nothing was omitted from any retained span, so the package carries no omission record.  No retained line contains a citation, so the wrapper carries no bibliography and no bibliography entry.  No retained line contains a figure environment, an image-inclusion command or drawing code, so no image is delivered and the package declares no asset dependency.  Editorial formatting only, in addition: an AMS \texttt{amsart} preamble that restates the article's own declarations and macros used by the retained lines, namely the environments \texttt{theorem} on separate counters, together with the standard packages those lines need.  The retained environments print in this document as Theorem 1 in order of appearance, whereas the article prints the same items with the numbering of its own full document.  The complete native source of the article as distributed in the arXiv e-print for this exact version is bundled unchanged as \texttt{sources/209766/source0.tex}.
\begin{equation}\label{eq:CGL}
  \begin{split}
    i \frac{\partial u}{\partial t} + (1 - i\epsilon)\Delta u + (1 + i\delta)|u|^{2}u &= 0 \qquad\text{in }\mathbb{R}^{3} \times (0, \infty),\\
    u|_{t=0} &= u_{0},
  \end{split}
\end{equation}

\begin{equation}\label{scaling2}
  u(x,t)\to\lambda^{1+i\varkappa}u(\lambda x,\lambda^{2}t),\qquad\lambda>0
\end{equation}

\begin{equation}\label{eq:perturbation-equation}
  \partial_\tau\phi=L_{\hat{Q}}\phi+\mathcal{N}(\phi),
\end{equation}

\begin{theorem}\label{thm:decaying-solution}
  Suppose that the operator \(L_{\hat{Q}}\) has at least one
  eigenvalue with positive real part. Then there exists \(\beta > 0\)
  such that for any \(r > 0\), we can find a mild solution
  \(\phi \in C((-\infty, 0], H^{2}(\mathbb{R}^{3}))\) to
  Equation~\eqref{eq:perturbation-equation} with
  \begin{equation*}
    0 < \sup_{\tau \in (-\infty, 0]} e^{-\beta\tau}\|\phi(\cdot, \tau)\|_{H^{2}(\mathbb{R}^{3})} < r.
  \end{equation*}
  Moreover, \(\phi\) can be taken to be real valued and such that
  \(\phi(\cdot, \tau) \neq 0\) for every \(\tau \leq 0\), and if the
  eigenfunction associated with the eigenvalue is radial, then
  \(\phi\) can be taken to be radial.
\end{theorem}

\begin{theorem}[Perturbed solution]\label{thm:perturbed-solution}
  Suppose that the assumptions of Theorem~\ref{thm:decaying-solution}
  hold and let \(\phi\) be a real-valued solution as in that theorem.
  With
  \(T = \frac{1}{2\kappa}\), \(\bm{\xi} = \frac{x}{\sqrt{2\kappa t}}\)
  and \(\tau = \log(2\kappa t)\), define
  \begin{align*}
    u_{\hat{Q}}(x, t) &= \frac{1}{(2\kappa t)^{\frac{1}{2}\left(1 + i\frac{\omega}{\kappa}\right)}}\hat{Q}(\bm{\xi}),\\
    u_{\phi}(x, t) &= \frac{1}{(2\kappa t)^{\frac{1}{2}\left(1 + i\frac{\omega}{\kappa}\right)}}\left(\hat{Q}(\bm{\xi}) + \phi(\bm{\xi}, \tau)\right),
  \end{align*}
  for \(t \in (0, T]\). Then, for every \(\alpha \in [0, 1 / 2)\), the
  difference \(u_{\phi} - u_{\hat{Q}}\) converges to zero in
  \(H^{\alpha}(\mathbb{R}^{3})\) as \(t \to 0^{+}\) and, extended by
  \(0\) at \(t = 0\), it satisfies
  \(u_{\phi} - u_{\hat{Q}} \in C([0, T], H^{\alpha}(\mathbb{R}^{3}))\).
\end{theorem}

  \begin{equation}\label{eq:energy-inequality}
    \int_{\mathbb{R}^{3}} |u(\cdot, t)|^{2}\psi(\cdot, t)
    + 2\epsilon \int_{-T}^{t}\int_{\mathbb{R}^{3}} |\nabla u|^{2}\psi
    \leq 2\real\left((\epsilon + i)\int_{-T}^{t}\int_{\mathbb{R}^{3}} u\nabla \conj{u}\nabla\psi\right)
    + \int_{-T}^{t}\int_{\mathbb{R}^{3}} |u|^{2}(\psi_{t} + 2\epsilon \Delta\psi).
  \end{equation}

\begin{equation}\label{eq:G}
  G(\mu, \gamma, \kappa) = \left(
    Q_{0}(\mu, \kappa; \xi_{1}) - Q_{\infty}(\gamma, \kappa; \xi_{1}),
    Q_{0}'(\mu, \kappa; \xi_{1}) - Q_{\infty}'(\gamma, \kappa; \xi_{1})
  \right) \colon \mathbb{R} \times \mathbb{C} \times \mathbb{R} \to \mathbb{C}^{2}.
\end{equation}

\begin{theorem}[Existence of a backward self-similar solution]\label{thm:backward-existence}
  For \(\epsilon \in \codenumber{[0.1681 \pm 10^{-15}]}\), the CGL
  equation has a backward self-similar solution. It is associated with
  a zero of the function \(G(\mu, \gamma, \kappa)\)~\eqref{eq:G},
  using \(\xi_{1} = \codenumber{16}\), with \((\mu, \gamma, \kappa)\)
  contained in
  \begin{align*}
    \mu &\in \resultnumber{2.2600773707_{82}^{95}},\\
    \gamma &\in \resultnumber{1.0571_{81}^{95} - 1.4043_{175}^{298} i},\\
    \kappa &\in \resultnumber{0.7912415632_{49}^{52}}.
  \end{align*}
  Moreover, as \(\xi \to \infty\), the profile satisfies
  \(Q(\xi) = p_{Q}\xi^{-1 - i\frac{\omega}{\kappa}} +
  \mathcal{O}(\xi^{-3})\), with
  \begin{equation*}
    p_{Q} \in \resultnumber{1.061_{84}^{93} - 1.400_{73}^{82} i}.
  \end{equation*}
\end{theorem}

\begin{theorem}\label{thm:main}
  The functions \(u_{1}\) and \(u_{2}\) both belong to
  \begin{equation*}
    L^{\infty}([-T, T], L^{3,\infty}(\mathbb{R}^{3}))
    \cap L_{\text{loc}}^{3}(\mathbb{R}^{3} \times (-T, T))
    \cap C^{\infty}((\mathbb{R}^{3} \times (-T, T)) \setminus \{(0, 0)\}),
  \end{equation*}
  are locally suitable weak solutions of~\eqref{eq:CGL} with the
  inequality in~\eqref{eq:energy-inequality} in fact being an
  equality, and share the same initial data
  \begin{equation*}
    u_{1}(\cdot, -T) = u_{2}(\cdot, -T) \in C^{\infty}(\mathbb{R}^{3}) \cap L^{3,\infty}(\mathbb{R}^{3}).
  \end{equation*}
  They satisfy
  \begin{equation*}
    u_{1} - u_{0}, u_{2} - u_{0} \in C([-T, T], H^{\alpha}(\mathbb{R}^{3}))\qquad \text{for any } \alpha \in [0, 1 / 2).
  \end{equation*}
  On \((0, T]\), only \(u_{1}\) is forward self-similar and they are
  distinct for all \(t \in (0, T]\).
\end{theorem}

\begin{proof}
  That \(u_{1}\) lies in
  \(L^{\infty}([-T, T], L^{3,\infty}(\mathbb{R}^{3}))\) follows from
  it having \(|x|^{-1}\) decay at infinity for all \(t\) and that the
  norm is invariant under the scaling~\eqref{scaling2}. For
  \(L_{\text{loc}}^{3}\) we use the substitution
  \(x = \sqrt{-2\kappa t}\,\bm{\xi}\), for which the \(L^{3}\) norm is
  invariant, together with \(|Q(\xi)| \lesssim \min(1, \xi^{-1})\), to
  get, for \(R \geq \sqrt{2\kappa T}\),
  \begin{equation*}
    \int_{|x| \leq R} |u_{Q}(x, t)|^{3}\,dx
    = \int_{|\bm{\xi}| \leq R / \sqrt{-2\kappa t}} |Q(|\bm{\xi}|)|^{3}\,d\bm{\xi}
    \lesssim 1 + \log\left(\frac{R}{\sqrt{-2\kappa t}}\right),
  \end{equation*}
  which is integrable in \(t\) on \([-T, 0)\). A similar computation
  for \(u_{\hat{Q}}\) gives us that
  \(u_{1} \in L_{\text{loc}}^{3}(\mathbb{R}^{3} \times (-T, T))\). For
  \(u_{2}\) the same arguments apply, with the difference
  \(w = u_{\phi} - u_{\hat{Q}}\) handled separately. Since the
  \(L^{3}(\mathbb{R}^{3})\) norm is invariant under the
  scaling~\eqref{scaling2}, we have
  \begin{equation*}
    \|w(\cdot, t)\|_{L^{3}(\mathbb{R}^{3})}
    = \|\phi(\cdot, \tau)\|_{L^{3}(\mathbb{R}^{3})}
    \lesssim \|\phi(\cdot, \tau)\|_{H^{2}(\mathbb{R}^{3})},
  \end{equation*}
  which by Theorem~\ref{thm:decaying-solution} is uniformly bounded
  for \(\tau \leq 0\). Together with
  \(L^{3}(\mathbb{R}^{3}) \subset L^{3,\infty}(\mathbb{R}^{3})\) this
  gives both the \(L^{3,\infty}\) and the \(L_{\text{loc}}^{3}\)
  bounds for \(u_{2}\).

  The smoothness for \(t \neq 0\) is immediate for \(u_{1}\) from the
  smoothness of \(Q\) and \(\hat{Q}\). The same applies to \(u_{2}\)
  for \(t < 0\), whereas for \(t > 0\) it is a consequence of a
  standard parabolic bootstrap. In particular, they both
  solve~\eqref{eq:CGL} pointwise for \(t \neq 0\). To prove smoothness
  across \(t = 0\) we first establish that both \(u_{1}\) and
  \(u_{2}\) are distributional solutions also across \(t = 0\).

  Note first that \(u_{1}, u_{2} \in L_{\text{loc}}^{3}\) gives
  \(|u_{1}|^{2}u_{1}, |u_{2}|^{2}u_{2} \in L_{\text{loc}}^{1}\), so
  that the distributional formulation makes sense. We next note that
  they are continuous across \(t = 0\) for \(x \neq 0\). For \(u_{1}\)
  this follows from \(Q\) and \(\hat{Q}\) having the same leading
  asymptotic behavior at infinity, so that both \(u_{Q}\) and
  \(u_{\hat{Q}}\) converge to \(u_{0}\) as \(t \to 0^{\mp}\), locally
  uniformly in \(x \neq 0\). For \(u_{2}\) we use that \(\phi\) is
  radial, so that the radial Sobolev inequality gives
  \(|\phi(\bm{\xi}, \tau)| \lesssim |\bm{\xi}|^{-1}\|\phi(\cdot,
  \tau)\|_{H^{1}(\mathbb{R}^{3})}\) and hence, the prefactor in \(w\)
  cancelling against \(|\bm{\xi}|^{-1}\),
  \begin{equation}\label{eq:w-pointwise-bound}
    |w(x, t)| \lesssim \frac{\|\phi(\cdot, \tau)\|_{H^{1}(\mathbb{R}^{3})}}{|x|}
    \lesssim \frac{(2\kappa t)^{\beta}}{|x|}
  \end{equation}
  by Theorem~\ref{thm:decaying-solution}. In particular, \(w\) is
  bounded near any point \((x_{0}, 0)\) with \(x_{0} \neq 0\) and
  tends to zero as \(t \to 0^{+}\), locally uniformly in
  \(x \neq 0\). Testing against a function supported away from the
  origin and splitting the integral into \(t < 0\) and \(t > 0\), the
  boundary terms at \(t = 0\) coming from the integration by parts in
  \(t\) cancel by this continuity, so that both solve~\eqref{eq:CGL}
  in \(\mathcal{D}'\) away from \((0, 0)\). It hence only remains to
  verify that the singularity at \((0, 0)\) is removable, which can be
  done directly, by multiplying the test function by a cutoff which
  cuts out a parabolic ball of radius \(\rho\) around the origin, and
  letting \(\rho \to 0\). They are hence solutions of~\eqref{eq:CGL}
  in \(\mathcal{D}'(\mathbb{R}^{3} \times (-T, T))\).

  This also gives us the smoothness across \(t = 0\). Indeed, both
  \(u_{1}\) and \(u_{2}\) are bounded near any point \((x_{0}, 0)\)
  with \(x_{0} \neq 0\), by~\eqref{eq:w-pointwise-bound} for the
  latter, and hence bounded solutions of~\eqref{eq:CGL} in
  \(\mathcal{D}'\) in a neighborhood of such a point. The same
  parabolic bootstrap as above therefore gives that they are smooth
  there, so that
  \(u_{1}, u_{2} \in C^{\infty}((\mathbb{R}^{3} \times (-T, T))
  \setminus \{(0, 0)\})\).

  It remains to verify that \(u_{1}\) and \(u_{2}\) are locally
  suitable weak solutions. The second condition is proved above, so we
  only have to check the first and the third one. For the first one,
  the bound \(|u_{1}| \lesssim \min(|t|^{-\frac{1}{2}}, |x|^{-1})\)
  gives
  \(u_{1} \in L^{\infty}([-T, T],
  L_{\text{loc}}^{2}(\mathbb{R}^{3}))\), and the substitution
  \(x = \sqrt{-2\kappa t}\,\bm{\xi}\) together with
  \(|Q'(\xi)| \lesssim \min(1, \xi^{-2})\) gives
  \begin{equation*}
    \int_{|x| \leq R} |\nabla u_{Q}(x, t)|^{2}\,dx
    = \frac{1}{\sqrt{-2\kappa t}}
    \int_{|\bm{\xi}| \leq R / \sqrt{-2\kappa t}} |Q'(|\bm{\xi}|)|^{2}\,d\bm{\xi}
    \lesssim \frac{1}{\sqrt{-2\kappa t}},
  \end{equation*}
  which is integrable in \(t\) on \([-T, 0)\). A similar computation
  for \(u_{\hat{Q}}\) gives us that
  \(u_{1} \in L^{2}([-T, T], W_{\text{loc}}^{1,2}(\mathbb{R}^{3}))\).
  For \(u_{2}\) the difference \(w\) is again handled separately: it
  is bounded in \(L^{3}(\mathbb{R}^{3})\), and hence in
  \(L_{\text{loc}}^{2}(\mathbb{R}^{3})\), uniformly in \(t\), and
  \begin{equation*}
    \|\nabla w(\cdot, t)\|_{L^{2}(\mathbb{R}^{3})}
    = (2\kappa t)^{-\frac{1}{4}}\|\nabla \phi(\cdot, \tau)\|_{L^{2}(\mathbb{R}^{3})}
    \lesssim (2\kappa t)^{\beta - \frac{1}{4}},
  \end{equation*}
  which is square integrable in \(t\) on \((0, T]\).

  For the third condition, multiplying~\eqref{eq:CGL} by
  \(\conj{u}\psi\), taking real parts and integrating by parts gives
  the identity corresponding to~\eqref{eq:energy-inequality}, with
  equality since \(\delta = 0\), for every \(\psi\) supported away
  from the origin. As for the distributional formulation, the general
  case follows by multiplying \(\psi\) with a cutoff which cuts out a
  parabolic ball of radius \(\rho\) around the origin, and letting
  \(\rho \to 0\).

  That \(u_{1}\) and \(u_{2}\) have the same initial data at
  \(t = -T\) is immediate from the definitions, both being given by
  \(u_{Q}\) for \(t < 0\). Since \(T = \frac{1}{2\kappa}\), it is
  given by \(Q(|\cdot|)\), which lies in
  \(C^{\infty}(\mathbb{R}^{3}) \cap L^{3,\infty}(\mathbb{R}^{3})\) by
  the above.

  For \(u_{1} - u_{0}\), let
  \(g(\bm{\xi}) = Q(|\bm{\xi}|) - p_{Q}|\bm{\xi}|^{-1 -
    i\frac{\omega}{\kappa}}\). Since \(Q\) is bounded, \(g\) behaves
  like \(|\bm{\xi}|^{-1}\) near the origin, and it is
  \(\mathcal{O}(|\bm{\xi}|^{-3})\) at infinity by
  Theorem~\ref{thm:backward-existence}, so that
  \(g \in H^{\alpha}(\mathbb{R}^{3})\) for every
  \(\alpha \in [0, 1/2)\). For \(t < 0\) we have
  \begin{equation*}
    u_{1}(x, t) - u_{0}(x)
    = \frac{1}{(-2\kappa t)^{\frac{1}{2}\left(1 + i\frac{\omega}{\kappa}\right)}}
    g\left(\frac{x}{\sqrt{-2\kappa t}}\right),
  \end{equation*}
  and hence, as in the proof of Theorem~\ref{thm:perturbed-solution},
  \(\|u_{1}(\cdot, t) - u_{0}\|_{H^{\alpha}(\mathbb{R}^{3})} \lesssim
  (-2\kappa t)^{\frac{1 -
      2\alpha}{4}}\|g\|_{H^{\alpha}(\mathbb{R}^{3})}\). The same holds
  for \(t > 0\) with \(\hat{Q}\) in place of \(Q\), and since
  \(u_{1}(\cdot, 0) = u_{0}\) this gives
  \(u_{1} - u_{0} \in C([-T, T], H^{\alpha}(\mathbb{R}^{3}))\). For
  \(u_{2}\) the same conclusion follows by combining this with
  Theorem~\ref{thm:perturbed-solution}, which gives
  \(w \in C([0, T], H^{\alpha}(\mathbb{R}^{3}))\) with
  \(w(\cdot, 0) = 0\).

  That \(u_{1}\) is forward self-similar on \((0, T]\) holds by
  construction. Self-similarity for \(u_{2}\) would require \(\phi\)
  to be independent of \(\tau\), which is impossible since \(\phi\) is
  non-zero but decays as \(\tau \to -\infty\) by
  Theorem~\ref{thm:decaying-solution}. Moreover, \(u_{1}\) and
  \(u_{2}\) are distinct for \(t \in (0, T]\) since \(\phi\) is
  non-zero for all \(\tau \leq 0\).
\end{proof}

\end{document}
